\section{Introduction}
\label{sec:introduction}

This document describes the winding connections of low-voltage distribution
transformers in two systems of practice. Section~\ref{sec:na} covers the
North American step from 480\,V to 240\,V inside a building: the single-phase
dual-voltage unit, the three-phase connections with a delta or wye secondary,
the open delta, and the checks that apply to all of them.
Section~\ref{sec:de} covers German practice, where the 400/230\,V system is
fed directly from the medium-voltage network: the vector groups of the
network transformer, the transformers used between two low-voltage levels,
the earthing systems, and a comparison with the North American connections.

Each connection is shown as a phasor diagram of its windings. A delta winding
is drawn as a triangle whose corners are the line terminals, a wye winding as
a star from the neutral point; the angle between the two diagrams is the
phase shift of the connection. Terminal markings follow the standard of each
system: H and X with X0 for the neutral in IEEE~C57.12.70~[1], and 1U/1V/1W,
2U/2V/2W with 2N for the neutral in IEC~60076-1~[2].

The values given are nominal system values and indicative ranges from the
standards cited. A project value follows from the equipment data and the
applicable edition of the code.
